% !TeX root = ../../tesis.tex

% -----------------------------------------------------------------------
\subsection{Coeficiente de Fricción Vial \texorpdfstring{($CF_h$)}{(CFh)}}
\label{subsec:friccion-vial}
% -----------------------------------------------------------------------

\textbf{Fuente principal:} \citet{fortin2016innovative}

\subsubsection{Definición y Fundamentación}

El Coeficiente de Fricción Vial es un multiplicador dinámico que corrige la
sobreestimación de eficiencia que ocurre cuando se asigna a las rutas de
transporte de superficie un tiempo de viaje ideal, sin considerar el efecto de
la congestión vehicular. \citet{fortin2016innovative} reconocen que los datos
estandarizados del \gtfs{} representan únicamente horarios planificados, y que
utilizar tiempos ajustados a las condiciones reales de tráfico devuelve
estimaciones sustancialmente más precisas al ejecutar el algoritmo de
\dijkstra{}.

Este indicador distingue dos categorías de modos con comportamientos
radicalmente distintos respecto a la congestión: los \termino{modos confinados}
(Metro, Cablebús, Suburbano, Tren Ligero), que circulan en infraestructura
segregada y mantienen un $CF = 1.0$ constante, y los
\termino{modos de tráfico mixto} (RTP, Trolebús, Metrobús en tramos compartidos,
microbuses concesionados), cuyo tiempo de viaje real se incrementa en función de
la hora del día y las condiciones de la vialidad.

\subsubsection{Fórmula}

\begin{equation}
    CF_h = \frac{T_{\text{real},h}}{T_{\text{ideal}}}
    \label{eq:friccion-vial}
\end{equation}

\begin{equation}
    T_{\text{real},h} = T_{\text{ideal}} \times CF_h
    \label{eq:tiempo-real}
\end{equation}

Donde:
\begin{itemize}
    \item $CF_h$ = Coeficiente de Fricción Vial en la franja horaria $h$
    (adimensional, $CF_h \geq 1.0$).
\item $T_{\text{real},h}$ = Tiempo de viaje observado en el tramo a la hora $h$,
    en minutos.
    \item $T_{\text{ideal}}$ = Tiempo de viaje en condiciones de flujo libre
    (referencia: 05:00\,h, sin congestión), en minutos.
\end{itemize}

\subsubsection{Ejemplo Práctico: \cdmx{} y Zona Metropolitana}

Los valores de $T_{\text{real},h}$ que permiten calcular $CF_h$ se obtienen del
\textbf{TomTom Traffic Index} \citep{tomtom2024}, base de datos histórica que
registra la velocidad promedio de cada corredor vial de la \cdmx{} por franja
horaria y día de la semana. Para cada arista de tráfico mixto del grafo, el
coeficiente se estima como el cociente entre el tiempo real medido a la hora $h$
y el tiempo de referencia en condiciones de flujo libre (05:00\,h).

\textbf{Tramo de referencia:} Ruta RTP, Insurgentes Norte (5\,km, tiempo en
flujo libre: 10 minutos).

\begin{table}[H]
    \centering
    \caption{Valores del coeficiente $CF_h$ por franja horaria.}
    \label{tab:friccion-vial}
    \begin{tabularx}{\textwidth}{X r r}
        \toprule
        \textbf{Franja horaria} & \textbf{$T_{\text{real}}$ observado} &
        \textbf{$CF_h$} \\
        \midrule
        05:00\,h (flujo libre)          & 10 min   & 1.00 (referencia) \\
        08:00\,h (pico máximo)          & 18.5 min & \textbf{1.85} \\
        15:00\,h (pico vespertino)      & 16 min   & 1.60 \\
        19:00\,h (pico nocturno)        & 17 min   & 1.70 \\
        Metro (cualquier hora)          & 10 min   & \textbf{1.00} (confinado) \\
        \bottomrule
    \end{tabularx}
\end{table}

El coeficiente $CF_h$ calibra las aristas de tráfico mixto del grafo construido
en el Objetivo Específico 1 (§\ref{subsec:obj-especificos}): sin esta
corrección, los datos \gtfs{} estáticos producirían un modelo que sobreestima la
eficiencia de los modos de superficie y subestima la brecha operativa entre los
modos confinados y los concesionados, distorsionando el diagnóstico de la red
actual que el Objetivo Específico 2 requiere.
